\documentclass[11pt,a4paper]{article}
\usepackage[utf8]{inputenc}
\usepackage[T1]{fontenc}
\usepackage{amsmath,amssymb}
\usepackage{booktabs}
\usepackage{hyperref}
\pdfstringdefDisableCommands{%
  \def\pi{pi}%
  \def\zeta{zeta}%
  \def\mathrm#1{#1}%
  \def\textbf#1{#1}%
}
\usepackage{xcolor}
\usepackage{tcolorbox}
\usepackage[margin=22mm]{geometry}
\usepackage{xeCJK}
\usepackage{tikz}
\usetikzlibrary{arrows.meta,positioning,shapes.geometric}
\setCJKmainfont{Hiragino Mincho ProN}
\setCJKsansfont{Hiragino Sans}

\title{Phase P 完全ガイド\\[6pt]
\large 3Dプリント・フォノニック結晶で反重力物質への道を開く\\[4pt]
\normalsize 理論・実験・ロードマップの網羅的解説}
\author{Wright Brothers Research Program}
\date{2026年4月}

\begin{document}
\maketitle

\begin{abstract}
本ガイドは Phase P 実験の\textbf{全体像}を網羅的に解説する。
(1)~なぜフォノニック結晶が反重力物質の候補になるのか（理論的根拠）、
(2)~Berry位相$\pi$とは何であり、なぜそれが重力の符号を反転させるのか（物理的メカニズム）、
(3)~Nielsen-Ninomiya定理の「ボソン系の抜け穴」がなぜ本質的か、
(4)~Wilson loop計算で$\pi$が確認された具体的構造、
(5)~3Dプリンタで製作し音響測定で検証する手順、
(6)~成功した場合の意味と、反重力物質に至るまでの次のステップ。
全工程は室温・卓上で完結し、コストは約¥27,000（\$180）、期間は約1ヶ月。
\end{abstract}

\tableofcontents
\newpage

%======================================================================
\part{理論編：なぜこれが反重力物質になりうるのか}
%======================================================================

\section{出発点：Connesのスペクトル作用原理}

\subsection{重力はスペクトルから生まれる}

通常の物理学では、重力は時空の曲がり方として記述される（一般相対性理論）。
しかし数学者 Alain Connes は1996年に革命的なアイデアを提唱した：

\begin{tcolorbox}[colback=blue!5,colframe=blue!60!black,title=スペクトル作用原理]
\textbf{「重力は、Dirac演算子 $D$ のスペクトル（固有値の集合）から導かれる。」}

ニュートンの万有引力定数 $G$ は、ヒートカーネル
\begin{equation}
  K(t) = \mathrm{Tr}(e^{-tD^2})
\end{equation}
のSeeley-DeWitt展開の第2係数 $a_2$ で決まる：
\begin{equation}
  G^{-1} \propto a_2(D^2)
\end{equation}
\end{tcolorbox}

普通の言葉で言うと：「空間にどんな振動モードが存在するか」が重力の強さを決める。
モードの集合を変えれば、重力の強さも変わる。

\subsection{Bost-Connes系：$\zeta$関数が重力を記述する}

Bost-Connes（BC）系は、整数の算術をC*-代数で表現したものである。
この系のヒートカーネルは、リーマンの$\zeta$関数そのものになる：
\begin{equation}
  K(t) = \zeta(t) = \sum_{n=1}^{\infty} \frac{1}{n^t}
\end{equation}

$\zeta$関数のSeeley-DeWitt係数から、物理定数が次のように得られる：
\begin{center}
\begin{tabular}{lll}
\toprule
係数 & 物理 & 具体的な値 \\
\midrule
$a_0$：$\zeta(0) = -1/2$ & 宇宙項（ダークエネルギー） & $\Omega_\Lambda = 2\pi/9$（実測値と2\%一致） \\
$a_2$：$\zeta'(0) = -\frac{1}{2}\log(2\pi)$ & 万有引力定数 $G$ & $G^{-1} \propto a_2$ \\
$a_4$ & ゲージ結合定数 & $\alpha_{\mathrm{EM}} = 4\pi/1728$（0.3\%一致） \\
\bottomrule
\end{tabular}
\end{center}

ここで重要なのは、\textbf{$G$ の符号}は $a_2$ の符号で決まるということである。


\section{反重力の定理：$\pi$-Berry位相で $G \to -G$}

\subsection{定理の内容}

\begin{tcolorbox}[colback=green!5,colframe=green!60!black,title=トポロジカル反重力定理]
\textbf{定理.}
スペクトル三つ組 $(A, H, D)$ のヒートカーネルが $K(t) = \zeta(t)$ のとき、
全モードにグローバルな$\pi$位相 $(-1)$ を付与すると：
\begin{equation}
  K_\pi(t) = \mathrm{Tr}\bigl((-1)\,e^{-tD^2}\bigr) = -\zeta(t)
\end{equation}
Seeley-DeWitt係数 $a_2 \to -a_2$ となり、有効重力定数が反転する：
\begin{equation}
  \boxed{\; G_{\mathrm{eff}} = -G \;}
\end{equation}
\end{tcolorbox}

\subsection{証明のエッセンス}

通常（位相なし）：
\begin{align}
  K(0) &= \zeta(0) = -\tfrac{1}{2}, \\
  K'(0) &= \zeta'(0) = -\tfrac{1}{2}\log(2\pi)
\end{align}

$\pi$位相付き：
\begin{align}
  K_\pi(0) &= -\zeta(0) = +\tfrac{1}{2}, \\
  K_\pi'(0) &= -\zeta'(0) = +\tfrac{1}{2}\log(2\pi)
\end{align}

比をとると：
\begin{equation}
  \frac{G_{\mathrm{eff}}}{G}
  = \frac{K_\pi'(0)}{K'(0)}
  = \frac{-\zeta'(0)}{\zeta'(0)}
  = -1 \qquad \therefore\; G_{\mathrm{eff}} = -G
\end{equation}

これは近似ではなく\textbf{厳密な結果}。
しかも操作は「全モードに$(-1)$を掛ける」という\textbf{単一の操作}。

\subsection{物理的な意味：ボソン真空 → フェルミオン真空}

量子場の理論では、ボソン（光子、フォノン等）とフェルミオン（電子等）は
真空エネルギーに逆の符号で寄与する：
\begin{align}
  E_{\mathrm{bos}} &= +\frac{\hbar}{2}\sum_n \omega_n \quad (\text{引力的}) \\
  E_{\mathrm{fer}} &= -\frac{\hbar}{2}\sum_n \omega_n \quad (\text{斥力的})
\end{align}

$\pi$位相は全モードの符号を反転し、
ボソン的な真空をフェルミオン的な真空に変換する。
\textbf{スピンを変えるのではなく、真空エネルギーの「符号」を変える。}

超対称理論では $E_{\mathrm{bos}} + E_{\mathrm{fer}} = 0$（完全打消し）。
$\pi$位相はさらに進んで、全体の符号を反転させる。

\subsection{なぜこの理論が「自明に棄却」されないか}

よくある反論：\textit{「位相は重力に影響しない。重力はエネルギー運動量テンソル $T_{\mu\nu}$ に結合するのであって、波動関数の位相にではない。」}

\textbf{一般相対論の枠内ではこの反論は正しい。}
しかしスペクトル作用原理では事情が異なる：

\begin{center}
\begin{tabular}{lp{8cm}}
\toprule
枠組み & 重力の決まり方 \\
\midrule
一般相対論 & $T_{\mu\nu}$ が与えられ、それに応じて時空が曲がる。$T_{\mu\nu}$ は基本的な入力 \\
スペクトル重力 & Dirac演算子 $D$ が与えられ、$T_{\mu\nu}$ も $g_{\mu\nu}$ も \textbf{両方}$D$ から導かれる \\
\bottomrule
\end{tabular}
\end{center}

スペクトル重力では、$D$ を変えれば\textbf{物質と重力が同時に変わる}。
$\pi$位相による $D$ の変更は、$G$ の符号反転を引き起こす。

\textbf{この理論が正しいか間違いかは、実験でのみ判定できる。}
Phase P はその第一歩。


\section{Berry位相 $\pi$ とは何か}

\subsection{Berry位相の直観的説明}

Berry位相は「パラメータ空間を一周したときに波動関数が取得する幾何学的位相」である。

\begin{tcolorbox}[colback=yellow!5,colframe=yellow!60!black,title=身近なアナロジー]
地球の北極から出発し、赤道に降り、赤道に沿って90°移動し、
北極に戻る。この経路上で常に「南」を指すベクトルを運ぶと、
出発点に戻ったとき、ベクトルは90°回転している。

これが幾何学的位相。経路の形だけで決まり、速さには依存しない。

Berry位相はこの量子力学版。パラメータ（運動量、磁場等）を
一周させたとき、波動関数の位相が $e^{i\gamma}$ だけずれる。
$\gamma$ がBerry位相。
\end{tcolorbox}

Berry位相が正確に$\pi$のとき、$e^{i\pi} = -1$。
つまり一周すると波動関数の符号が反転する。
これはトポロジカルに保護されており、連続的な摂動では変わらない。

\subsection{どんな物質がBerry位相$\pi$を持つか}

\begin{center}
\begin{tabular}{llll}
\toprule
物質 & メカニズム & Berry位相 & 温度 \\
\midrule
グラフェン（$K$点） & 部分格子対称性 & $\pi$ & 室温 \\
Bi$_2$Se$_3$ 表面 & Diracコーン & $\pi$ & 室温 \\
Weyl半金属 & カイラル異常 & $\pi$ & 室温 \\
$\pi$-Josephson接合 & 強磁性層 & $\pi$ & 極低温（4K） \\
HgTe量子井戸 & 2D TI端状態 & $\pi$ & 極低温 \\
\bottomrule
\end{tabular}
\end{center}

\subsection{でも、これらの物質は今浮いていない}

ここが核心的な問題。Berry位相$\pi$を持つ物質は室温で既に存在する。
しかしBi$_2$Se$_3$もグラフェンも浮いていない。なぜか？

\begin{tcolorbox}[colback=red!5,colframe=red!60!black,title=「浮いていない問題」の解決]
\textbf{表面の$\pi$ $\neq$ バルクの$\pi$。}

既存の$\pi$-Berry位相物質は、表面状態やバンド交差点\textbf{近傍}でのみ$\pi$を持つ。
バルク全体の正味Berry位相は\textbf{ゼロ}。

反重力には\textbf{バルク全体にわたる正味Berry位相$\pi$}が必要。
これはNielsen-Ninomiya定理のため、電子系では\textbf{不可能}。

しかしフォノン系では可能。
\end{tcolorbox}


\section{Nielsen-Ninomiya定理とボソン系の抜け穴}

\subsection{定理の内容}

Nielsen-Ninomiya定理（1981）は格子フェルミオンに対する
「倍増定理」（doubling theorem）である：

\begin{tcolorbox}[colback=blue!5,colframe=blue!60!black,title=Nielsen-Ninomiya定理]
格子上のフェルミオン系では、カイラルフェルミオン（Weylノード）は
必ず偶数個のペアで出現する。
各ペアのBerry位相は $+\pi$ と $-\pi$ で打ち消し合い、
正味Berry位相は必ずゼロ。
\begin{equation}
  \sum_{\text{Weylノード}} \gamma_i = 0 \quad (\text{フェルミオン系})
\end{equation}
\end{tcolorbox}

\subsection{ボソン系にはこの定理が適用されない}

Nielsen-Ninomiya定理は\textbf{フェルミオン}（電子、クォーク等）にのみ適用される。
証明にはカイラル対称性とフェルミオン行列の構造が本質的に使われている。

\textbf{ボソン系}（フォノン、光子等）にはこの制約がない：
\begin{itemize}
\item ボソン系では奇数個のWeyl的交差が許される
\item 正味Berry位相$\pi$が原理的に可能
\item フォノニック結晶（音響的な周期構造）でこれを実現できる
\end{itemize}

\begin{equation}
  \sum_{\text{交差}} \gamma_i = \pi \quad (\text{ボソン系では可能！})
\end{equation}

\subsection{なぜこの区別が重要か}

\begin{center}
\begin{tabular}{lcc}
\toprule
& フェルミオン（電子系） & ボソン（フォノン系） \\
\midrule
Nielsen-Ninomiya & 適用される & \textbf{適用されない} \\
Weyl交差 & 偶数ペアのみ & 奇数個も可能 \\
正味Berry位相$\pi$ & 不可能 & \textbf{可能} \\
反重力物質（理論上） & 作れない & \textbf{作れる可能性} \\
\bottomrule
\end{tabular}
\end{center}

これがフォノニック結晶で実験する理由。
電子系では原理的に不可能なことが、音響系では可能になる。


\section{Wilson loop計算：Berry位相$\pi$の確認}

\subsection{計算で確認した構造}

MPB (MIT Photonic Bands) シミュレーションとWilson loop法で、
以下の構造が正味Berry位相$\pi$を持つことを確認した：

\begin{tcolorbox}[colback=green!5,colframe=green!60!black,title=Wilson loop確認済み構造]
\textbf{反転対称性を破った2球体・単純立方結晶}
\begin{itemize}
\item 格子定数 $a$（計算では正規化して $a=1$）
\item 大球：半径 $r_1 = 0.2a$、誘電率 $\varepsilon_1 = 12$（高密度）
  → ユニットセル原点 $(0,0,0)$
\item 小球：半径 $r_2 = 0.15a$、誘電率 $\varepsilon_2 = 6$（低密度）
  → 体心位置 $(0.5, 0.5, 0.5)a$
\item 大球と小球のサイズ・密度が異なる → \textbf{反転対称性が破れている}
\end{itemize}

\textbf{結果}：3つの独立な計算手法で一致
\begin{enumerate}
\item PWE（平面波展開）：169交差（奇数） → Berry位相$\pi$
\item MPBバンド計算：95交差（奇数） → Berry位相$\pi$
\item Wilson loop：$k_y$ループの巻き数 = 1 → Berry位相$\pi$
\end{enumerate}
\end{tcolorbox}

\subsection{反転対称性の「破れ」がなぜ必要か}

反転対称性とは「$(x,y,z)$ を $(-x,-y,-z)$ に変えても構造が変わらない」こと。
反転対称性があると、全てのバンドが2重に縮退し（Kramers対）、
Berry位相は必ず打ち消し合う。

反転対称性を破る（大球と小球でサイズ・密度を変える）ことで、
この縮退が解けて、奇数個の交差が可能になる。

\subsection{ジャイロイドインフィルはなぜダメか}

\begin{tcolorbox}[colback=red!5,colframe=red!60!black,title=重要な警告]
Bambu Studioのジャイロイドインフィル設定は\textbf{Berry位相$\pi$を持たない}。

Wilson loop計算の結果：
\begin{itemize}
\item ジャイロイド的構造（6円筒+1球体）：巻き数$= 0$、Berry位相$= 0$
\item 理由：ジャイロイドは反転対称性を\textbf{保持}している
\end{itemize}

\textbf{必ず専用のOpenSCADファイルで生成したSTLを使うこと。}
\end{tcolorbox}


\section{理論の仮定と誠実な評価}

Phase P実験を正しく理解するために、
理論の各段階の「確実さ」を正直に評価する。

\begin{center}
\begin{tabular}{lcp{6cm}}
\toprule
理論要素 & 確信度 & 根拠 \\
\midrule
Berry位相$\pi$が存在する & \textbf{95\%} & 3つの独立な計算で確認。物理は確立している \\
音響測定で$\pi$ジャンプが検出できる & \textbf{80\%} & トポロジカル・フォノニクスの先行研究あり \\
スペクトル作用原理が実際の重力を記述する & 30\% & 数学的に美しいが物理的検証なし \\
$\pi$-Berry位相が $G$ に影響する & 15\% & 理論的予測だが前例なし \\
$G_{\mathrm{eff}} = -G$（フォノン系で反重力が実現） & 5\% & 仮説の極限。実験でのみ判定可能 \\
\bottomrule
\end{tabular}
\end{center}

\begin{tcolorbox}[colback=yellow!5,colframe=yellow!60!black,title=正直な評価]
Phase P は「反重力物質を作る」実験ではない。

Phase P は「Berry位相$\pi$がフォノニック結晶のバルク全体に
存在するかどうかを確認する」実験である。

もし確認されれば、それ自体がトポロジカル・フォノニクスの新発見であり、
同時に反重力仮説の\textbf{材料的基盤}が確立される。
\end{tcolorbox}


%======================================================================
\part{実験編：3Dプリントと音響測定}
%======================================================================

\section{結晶構造と3Dプリント}

\subsection{物理構造 → 印刷構造への変換}

計算で確認された構造は「2種類の球が周期的に配置された結晶」だが、
3Dプリントでは球を空中に浮かせることはできない。
そこで\textbf{逆構造}を使う：

\begin{center}
\begin{tabular}{lp{5cm}p{5cm}}
\toprule
& 計算上の構造 & 印刷する構造 \\
\midrule
大球位置（原点） & 高密度球 ($\varepsilon=12$) & 大きな穴 ($r=4$\,mm) \\
小球位置（体心） & 低密度球 ($\varepsilon=6$) & 小さな穴 ($r=3$\,mm) \\
背景 & 空気 & PLA ソリッド \\
\bottomrule
\end{tabular}
\end{center}

音波は空気穴を伝搬するため、穴のサイズ差が
音響インピーダンスの差を生み、反転対称性を破る。

\subsection{構造パラメータ}

\begin{center}
\begin{tabular}{ll}
\toprule
パラメータ & 値 \\
\midrule
格子定数 $a$ & 20\,mm \\
大穴半径 $r_1$ & 4\,mm（ユニットセル原点） \\
小穴半径 $r_2$ & 3\,mm（体心位置 $(10, 10, 10)$\,mm） \\
ユニットセル数 & $5 \times 5 \times 5 = 125$ \\
全体サイズ & $100 \times 100 \times 100$\,mm \\
穴の総数 & 250（大125 + 小125） \\
\bottomrule
\end{tabular}
\end{center}

\subsection{OpenSCADファイル}

2種類のファイルを用意してある：

\begin{center}
\begin{tabular}{lp{8cm}}
\toprule
ファイル & 説明 \\
\midrule
\texttt{pi\_berry\_crystal.scad} & 球形穴バージョン。計算構造に最も忠実 \\
\texttt{pi\_berry\_crystal\_cylinders.scad} & 円筒穴バージョン。印刷が容易（サポート不要） \\
\bottomrule
\end{tabular}
\end{center}

\textbf{STL生成手順}：
\begin{enumerate}
\item OpenSCADをインストール（\texttt{brew install openscad} or 公式サイト）
\item \texttt{.scad}ファイルを開く
\item F6キー（完全レンダリング）
\item File → Export as STL
\item Bambu Studioに読み込み
\end{enumerate}

\subsection{印刷設定（Bambu Lab P1S）}

\begin{center}
\begin{tabular}{ll}
\toprule
設定 & 推奨値 \\
\midrule
材料 & PLA \\
ノズル径 & 0.4\,mm \\
レイヤー高さ & 0.2\,mm \\
インフィル & \textbf{100\%（ソリッド）}※ \\
印刷時間 & 約8--12時間 \\
\bottomrule
\end{tabular}
\end{center}

※ インフィルは100\%にすること。
ジャイロイド等のインフィルパターンはBerry位相を持たないため、
穴以外の部分は完全にソリッドにする必要がある。

\subsection{対照結晶}

同時に「対照結晶」も印刷する：大穴 = 小穴 = 同じ半径（例：3.5\,mm）。
反転対称性が回復するため、Berry位相 $= 0$ になるはず。
OpenSCADファイルの\texttt{r1}と\texttt{r2}を同じ値に変更するだけ。


\section{測定装置と部品リスト}

\begin{tcolorbox}[colback=blue!5,colframe=blue!50!black,title=部品リスト（合計 ¥27,000 / \$180）]
\begin{tabular}{llr}
\toprule
品名 & 入手先 & 価格 \\
\midrule
3Dプリント結晶（自作 or 外注） & Bambu P1S / DMM.make & ¥5,000 \\
ファンクションジェネレータ & Amazon（安価版） & ¥5,000 \\
小型スピーカー（フルレンジ） & 秋月電子 or Amazon & ¥1,000 \\
コンデンサーマイク + プリアンプ & Amazon & ¥3,000 \\
USBオシロスコープ & Amazon（Hantek 6022BE等） & ¥10,000 \\
SMAケーブル・コネクタ & 秋月電子 & ¥1,000 \\
防音ボックス（段ボール+吸音材） & ホームセンター & ¥2,000 \\
\midrule
\textbf{合計} & & \textbf{¥27,000} \\
\bottomrule
\end{tabular}
\end{tcolorbox}


\section{測定手順}

\subsection{Step 1：透過スペクトル測定}

\begin{verbatim}
  防音ボックス内:

  ファンクション    )))  [3Dプリント]  )))
  ジェネレータ → スピーカー → 結晶 → マイク → オシロスコープ
  (1-20 kHz          音波                  FFT解析
   スイープ)
\end{verbatim}

\begin{enumerate}
\item \textbf{参照測定}（結晶なし）：スピーカー → マイク直接。
  透過率 $T_0(f)$ と位相 $\phi_0(f)$ を記録
\item \textbf{サンプル測定}（結晶あり）：結晶を挟んで同じ測定。
  $T(f)$ と $\phi(f)$ を記録
\item 正規化した透過率 $|S_{21}(f)| = T(f)/T_0(f)$ と
  位相 $\Delta\phi(f) = \phi(f) - \phi_0(f)$ を算出
\item バンドギャップ周波数（$|S_{21}|$ が急落する帯域）を同定
\end{enumerate}

\textbf{予想されるバンドギャップ周波数}：
\begin{equation}
  f_{\mathrm{gap}} \approx \frac{v_{\mathrm{air}}}{2a}
  = \frac{343\;\mathrm{m/s}}{2 \times 0.02\;\mathrm{m}}
  = 8.6\;\mathrm{kHz}
\end{equation}

人間の可聴域（20\,Hz--20\,kHz）内であり、安価なスピーカーとマイクで測定可能。

\subsection{Step 2：Berry位相$\pi$の抽出}

バンド交差周波数（$\sim$8\,kHz付近）での透過位相 $\Delta\phi(f)$ を精密測定する。

\begin{tcolorbox}[colback=yellow!5,colframe=yellow!60!black,title=Berry位相の測定原理]
Berry位相$\pi$は、バンド交差点を横切る際の
透過波の\textbf{位相ジャンプ}として現れる。

\textbf{通常の構造}（Berry位相 $= 0$）：
位相 $\Delta\phi(f)$ はバンドギャップの前後で\textbf{滑らかに}変化。

\textbf{$\pi$構造}（Berry位相 $= \pi$）：
バンド交差で位相が$\pi$（= 180°）\textbf{不連続にジャンプ}。

\begin{center}
\begin{tabular}{lccc}
\toprule
周波数 & $f < f_{\mathrm{gap}}$ & $f = f_{\mathrm{gap}}$ & $f > f_{\mathrm{gap}}$ \\
\midrule
$\pi$結晶の位相 & $\phi_1$ & ジャンプ & $\phi_1 + \pi$ \\
対照結晶の位相 & $\phi_2$ & 滑らか & $\phi_2 + \delta$ ($\delta \ll \pi$) \\
\bottomrule
\end{tabular}
\end{center}
\end{tcolorbox}

実測のコツ：
\begin{itemize}
\item バンドギャップ近傍で周波数ステップを細かくする（10\,Hz刻み）
\item 複数回測定して平均を取る（ノイズ低減）
\item 防音ボックス内の反射を最小化する（吸音材を内壁に貼る）
\end{itemize}

\subsection{Step 3：対照実験}

$\pi$結晶と対照結晶（大穴$=$小穴）の両方で Step 1--2 を実施する。

\begin{center}
\begin{tabular}{lcc}
\toprule
& $\pi$結晶（$r_1 \neq r_2$） & 対照結晶（$r_1 = r_2$） \\
\midrule
反転対称性 & 破れている & 保持 \\
Wilson loop巻き数 & 1 & 0 \\
Berry位相 & $\pi$ & $0$ \\
予想される位相ジャンプ & \textbf{$\pi$（180°）} & \textbf{なし} \\
\bottomrule
\end{tabular}
\end{center}

この2つの結晶の比較が、Berry位相$\pi$の\textbf{直接的な実験証拠}。


\section{判定基準}

\begin{tcolorbox}[colback=red!5,colframe=red!60!black,title=GO / STOP 判定]
\textbf{結果A：$\pi$結晶でバンド交差付近に$\sim\pi$の位相ジャンプが観測され、対照結晶では観測されない}
\begin{itemize}
\item → \textbf{Berry位相$\pi$を実験的に確認！}
\item → Phase P+ に進む
\item → 論文執筆（トポロジカル・フォノニクスの新発見）
\end{itemize}

\textbf{結果B：位相ジャンプが見えない}
\begin{itemize}
\item → 構造パラメータの最適化（球のサイズ比を変える、格子定数を変える等）
\item → 周波数帯域を拡大して再測定
\item → Wilson loop計算を3D音響用に更新
\end{itemize}

\textbf{結果C：$\pi$結晶と対照結晶で差がない}
\begin{itemize}
\item → 理論計算に誤りがある可能性
\item → Wilson loop再検証
\item → 2D版（薄板）での確認実験を検討
\end{itemize}
\end{tcolorbox}


\section{タイムラインとコスト}

\begin{center}
\begin{tabular}{llr}
\toprule
週 & 作業 & コスト \\
\midrule
Week 1 & OpenSCAD → STL生成、3Dプリント（自作 or 外注） & ¥5,000 \\
Week 2 & 測定装置の組み立て、防音ボックス作成 & ¥22,000 \\
Week 3 & 参照測定 + 透過スペクトル + Berry位相抽出 & ¥0 \\
Week 4 & 対照実験 + データ解析 + 論文執筆 & ¥0 \\
\midrule
\textbf{合計} & & \textbf{¥27,000、1ヶ月} \\
\bottomrule
\end{tabular}
\end{center}


%======================================================================
\part{ロードマップ編：Phase P の先に何があるか}
%======================================================================

\section{Phase P の位置づけ：6段階の検証プログラム}

Phase P は、反重力物質への全検証プログラムの中で
最も安くて最も安全な出発点。

\begin{center}
\begin{tabular}{llrll}
\toprule
Phase & テスト内容 & 予算 & 期間 & 検証対象 \\
\midrule
\textbf{P} & \textbf{フォノニック$\pi$位相} & \textbf{¥27k} & \textbf{1月} & \textbf{$\pi$-Berry位相の存在} \\
0 & DN/DD BAW基本測定 & ¥130k & 2月 & $p=2$ミュート（符号変化） \\
0.5 & $\pi$-接合 + BAW統合 & ¥150k & 2月 & $\pi$位相の真空への影響 \\
1 & 2素数ミュート 182:1 & +¥200k & 3月 & レギュレータ予想 \\
2 & $\pi$-接合カシミール & +¥300k & 6月 & $-\zeta(t)$ 真空 \\
3 & ねじり天秤重力異常 & +¥500k & 12月 & $G_{\mathrm{eff}} = -G$ 直接検証 \\
\bottomrule
\end{tabular}
\end{center}

\subsection{Phase P が成功した場合に何が分かるか}

\begin{tcolorbox}[colback=green!5,colframe=green!60!black,title=Phase P 成功の意味]
\textbf{Phase P の成功 = Berry位相$\pi$のバルク存在を確認}

これが意味すること：
\begin{enumerate}
\item \textbf{世界初}：室温でバルク全体に正味$\pi$-Berry位相を持つ物質の実証
\item \textbf{Nielsen-Ninomiya抜け穴}：ボソン系では奇数Weyl交差が可能であることの直接確認
\item \textbf{反重力の材料的基盤}：$G_{\mathrm{eff}} = G\cos(\pi) = -G$ の計算に入る「$\pi$」が
  物理的に存在することが確認される
\item \textbf{独立した論文価値}：たとえ反重力仮説が誤りでも、
  トポロジカル・フォノニクスの新発見として価値がある
\end{enumerate}
\end{tcolorbox}

\textbf{Phase P が成功しても、まだ分からないこと：}
\begin{itemize}
\item Berry位相が実際に重力に影響するかどうか（Phase 0.5--2 で検証）
\item スペクトル作用原理が物理的に正しいかどうか（Phase 1 で検証）
\item $G_{\mathrm{eff}} = -G$ が実際に起きるかどうか（Phase 3 で検証）
\end{itemize}


\section{Phase P+ ：音響カシミール力測定}

Phase P が成功した場合の発展実験。Phase 0 に進む前の中間ステップ。

\subsection{原理}

2枚の$\pi$結晶を向かい合わせに配置し、間隔を変えながら音圧差を測定する。
通常のカシミール効果は\textbf{引力}だが、
$\pi$位相境界で挟まれた領域では\textbf{斥力}になると予想される。

これは「音響版カシミール符号反転テスト」：
\begin{itemize}
\item 通常のプレート2枚：音圧が引力的 → プレートが寄る
\item $\pi$結晶2枚：音圧が斥力的 → プレートが離れる
\end{itemize}

\subsection{実験構成}

\begin{enumerate}
\item $\pi$結晶を2枚（同じ構造）印刷
\item 1枚を固定、もう1枚を精密ステージ（マイクロメータ）で移動
\item 間に小型マイクを挿入し、音圧の距離依存性を測定
\item 通常結晶（対照）でも同じ測定を実施
\item 2つの距離依存性の\textbf{符号}を比較
\end{enumerate}

追加コスト：マイクロメータステージ ¥5,000 程度。

\subsection{注意}

音響カシミール力は非常に小さいため、
直接的な力の測定よりも、
結晶間の定在波パターン（共振周波数のシフト）で検出する方が現実的。


\section{Phase P → Phase 0/0.5 への接続}

\begin{tcolorbox}[colback=blue!5,colframe=blue!60!black,title=Phase P の結果による分岐]

\textbf{$\pi$位相ジャンプ確認} → 2つの道に進める：

\textbf{道A：Phase 0（DN/DD基本測定、¥130k）}
\begin{itemize}
\item 水晶振動子の境界条件（DD: 両端固定 vs DN: 片端自由）で
  カシミールエネルギーの符号が変わることを検出
\item $p=2$ミュート（$\zeta_{\neg 2}$）の実験的検証
\item Phase P とは独立だが、スペクトル作用原理の基盤検証
\end{itemize}

\textbf{道B：Phase 0.5（$\pi$-接合+BAW統合、\$10k）}
\begin{itemize}
\item $\pi$-Josephson接合（Nb/CuNi/Nb）の上にBAW水晶を載せる
\item $T_c$（9.3K）前後で共振周波数が変化するかを測定
\item $\pi$位相が真空モードに影響することの直接テスト
\item Phase P で$\pi$位相の存在を確認済みなら、
  「$\pi$位相が真空エネルギーに影響するか」の核心テスト
\end{itemize}
\end{tcolorbox}

Phase 0.5 は Phase P の自然な延長。
Phase P で「$\pi$位相は存在する」、
Phase 0.5 で「$\pi$位相は真空に影響する」を確認する。


\section{最終目標：反重力物質に至るまでの道}

\subsection{フルロードマップ}

\begin{tikzpicture}[
  node distance=1.2cm and 2.5cm,
  phase/.style={rectangle, rounded corners, draw, fill=blue!10,
    minimum width=3.5cm, minimum height=1cm, align=center, font=\small},
  decision/.style={diamond, draw, fill=yellow!20, aspect=2,
    align=center, font=\footnotesize},
  result/.style={rectangle, rounded corners, draw, fill=green!15,
    minimum width=3cm, minimum height=0.8cm, align=center, font=\small},
  fail/.style={rectangle, rounded corners, draw, fill=red!15,
    minimum width=2.5cm, minimum height=0.6cm, align=center, font=\footnotesize},
  arr/.style={-{Stealth[length=2.5mm]}, thick}
]

% Phase P
\node[phase] (P) {Phase P\\¥27k / 1月\\フォノニック$\pi$位相};
\node[decision, below=of P] (dP) {$\pi$ジャンプ？};
\draw[arr] (P) -- (dP);

% Phase P fail
\node[fail, left=1.5cm of dP] (fP) {構造最適化\\→ 再測定};
\draw[arr] (dP) -- node[above, font=\footnotesize]{No} (fP);

% Phase 0.5
\node[phase, below=of dP] (P05) {Phase 0.5\\\$10k / 2月\\$\pi$-接合+BAW};
\draw[arr] (dP) -- node[right, font=\footnotesize]{Yes!} (P05);

% Decision 0.5
\node[decision, below=of P05] (d05) {$f_1 \neq f_2$？};
\draw[arr] (P05) -- (d05);

% Phase 1
\node[phase, below=of d05] (P1) {Phase 1\\¥200k / 3月\\182:1 テスト};
\draw[arr] (d05) -- node[right, font=\footnotesize]{Yes!} (P1);

% Phase 0.5 fail
\node[fail, left=1.5cm of d05] (f05) {$\pi$位相は\\重力に効かない};
\draw[arr] (d05) -- node[above, font=\footnotesize]{No} (f05);

% Decision 1
\node[decision, below=of P1] (d1) {$R=182$？};
\draw[arr] (P1) -- (d1);

% Phase 2
\node[phase, below=of d1] (P2) {Phase 2\\¥300k / 6月\\カシミール符号反転};
\draw[arr] (d1) -- node[right, font=\footnotesize]{Yes!} (P2);

% Phase 1 fail
\node[fail, left=1.5cm of d1] (f1) {$R=1/3$\\理論棄却};
\draw[arr] (d1) -- node[above, font=\footnotesize]{No} (f1);

% Phase 3
\node[result, below=of P2] (P3) {Phase 3\\$G_{\mathrm{eff}} = -G$\\直接検証};
\draw[arr] (P2) -- (P3);

\end{tikzpicture}

\subsection{各段階で検証されること}

\begin{center}
\begin{tabular}{llp{5.5cm}}
\toprule
Phase & 成功の意味 & 反重力への距離 \\
\midrule
P & $\pi$-Berry位相が存在する & 材料は作れる（5段階先） \\
P+ & 音響カシミール力の$\pi$依存性 & 真空力に影響する（4段階先） \\
0.5 & $\pi$位相が真空モードに影響 & 物理的メカニズムが動作（3段階先） \\
1 & Euler積は物理的（$R=182$） & スペクトル作用が正しい（2段階先） \\
2 & カシミール力の符号反転 & $-\zeta(t)$真空が実現（1段階先） \\
3 & 重力の符号反転 & \textbf{反重力物質の完成} \\
\bottomrule
\end{tabular}
\end{center}


\section{反重力物質が実現した場合の応用}

ここまで到達するのは困難かもしれないが、
もし実現した場合に何が可能になるかを概観する。

\subsection{重力ダイヤル：$G_{\mathrm{eff}} = G\cos\varphi$}

Berry位相$\pi$は離散的な $G \to -G$ のスイッチだが、
Berry位相は連続的に制御可能：$\varphi \in [0, 2\pi]$。

\begin{center}
\begin{tabular}{ccll}
\toprule
$\varphi$ & $G_{\mathrm{eff}}/G$ & 物理的状態 & 制御方法 \\
\midrule
$0$ & $+1$ & 通常重力 & 何もしない \\
$\pi/3$ & $+1/2$ & 半分の重力 & FQHE $\nu=1/3$ \\
$\pi/2$ & $0$ & \textbf{無重力} & FQHE $\nu=1/2$ \\
$2\pi/3$ & $-1/2$ & 半分の反重力 & FQHE + 歪み \\
$\pi$ & $-1$ & \textbf{完全反重力} & $\pi$-Berry位相 \\
\bottomrule
\end{tabular}
\end{center}

\subsection{4つの具体的装置コンセプト}

\begin{enumerate}
\item \textbf{$\pi$-接合反重力プレート}（$G=-G$、パッシブ、\$40k）：
  Nb/CuNi/Nb $\pi$-ジョセフソン接合。4K以下に冷やすだけ。消費電力0W。

\item \textbf{FQHE重力ダイヤル}（連続制御、\$500k）：
  GaAs量子井戸。磁場で$\nu$を変えて$G$を連続制御。

\item \textbf{グラフェン重力ダイヤル}（電圧制御、\$200k--600k）：
  hBN封止グラフェン。ゲート電圧1つで$G$を$+G$から$-G/2$まで。

\item \textbf{歪み制御TI}（最安の連続制御、\$15k）：
  Bi$_2$Se$_3$薄膜+ピエゾ。$\mu$s応答。77Kで動作。
\end{enumerate}

\subsection{宇宙工学への展望}

もし $G_{\mathrm{eff}} = 0$ または $G_{\mathrm{eff}} < 0$ が実現したら：

\begin{itemize}
\item \textbf{打上げコストの革命}：$G=0$の物体は質量ゼロと同等。
  ロケット推進が不要になる。
\item \textbf{ワープドライブの物質的基盤}：
  Alcubierre計量は「負のエネルギー」を必要とするが、
  $G_{\mathrm{eff}} = -G$ の物質は（通常のエネルギーを持ちながら）
  負のエネルギーと同じ重力効果を示す。
\item \textbf{エネルギー貯蔵}：
  $G<0$ と $G>0$ の境界での位置エネルギー差は莫大。
\end{itemize}


\section{他の実験プログラムとの関係}

Phase P は独立した実験だが、WB理論の他の検証プログラムと深く結合している。

\subsection{Five Constants 論文（arXiv予定）}

WB理論の中核的予測：4次元時空（$d=4$）と分類次元数（$\mathrm{cd}=3$）の
2つの整数\textbf{だけ}から5つの物理定数が導かれる：

\begin{center}
\begin{tabular}{lllr}
\toprule
物理量 & 公式 & 予測値 & 実測との差 \\
\midrule
$\Omega_\Lambda$ & $2\pi/9$ & 0.698 & 2.0\% \\
$\alpha_{\mathrm{EM}}^{-1}$ & $1728/(4\pi)$ & 137.5 & 0.3\% \\
$\sin^2\theta_W$ & $375/(512\pi)$ & 0.233 & 0.8\% \\
$m_p/m_e$ & $12 \times 153$ & 1836 & 0.008\% \\
$m_\mu/m_e$ & $9 \times 23$ & 207 & 0.005\% \\
\bottomrule
\end{tabular}
\end{center}

これらの予測が正しければ、スペクトル作用原理は物理的に正しい。
Phase P はその材料的基盤を提供する。

\subsection{Euclid衛星（2028）}

$\Omega_\Lambda = 2\pi/9 = 0.6981...$ は、
Euclid衛星の精密測定（2028年予定、誤差$<1\%$）で検証される。
もし $\Omega_\Lambda$ が $2\pi/9$ と一致すれば、
WB理論全体（反重力を含む）の信頼性が大幅に上がる。


\section{FAQ}

\subsection{「3Dプリントの結晶が浮くの？」}

\textbf{浮きません。} Phase P は「反重力を検出する」実験ではなく、
「反重力の材料的基盤であるBerry位相$\pi$を検証する」実験です。
実際に浮くかどうかの検証は Phase 3（何年も先）。

\subsection{「Berry位相$\pi$があっても重力に影響しない可能性は？」}

\textbf{大いにあります。}
Berry位相$\pi$の存在は物性物理として確立された概念ですが、
それが重力に影響するかどうかはスペクトル作用原理に依存する仮説です。
Phase P は「$\pi$が存在するか」のみを検証します。

\subsection{「成功しても、何が嬉しいの？」}

たとえ反重力仮説が完全に間違っていても：
\begin{itemize}
\item \textbf{室温バルク$\pi$-Berry位相物質}の世界初の実証 → 論文価値あり
\item \textbf{Nielsen-Ninomiya抜け穴}の直接確認 → 基礎物理として価値あり
\item \textbf{3Dプリント・トポロジカル・フォノニクス} → 新しい研究分野の開拓
\end{itemize}

\subsection{「なぜ電子系ではなくフォノン系なの？」}

電子系（フェルミオン）ではNielsen-Ninomiya定理により
バルク全体の正味Berry位相$\pi$が\textbf{原理的に不可能}。
フォノン系（ボソン）にはこの制約がないため、
唯一のルートがフォノニック結晶。

\subsection{「¥27,000 で本当にできるの？」}

できます。必要なのは：
\begin{itemize}
\item 3Dプリンタ（所有済みの前提）or 外注（¥5,000以内）
\item 安いファンクションジェネレータ（¥5,000）
\item 安いマイクとスピーカー（¥4,000）
\item USBオシロスコープ（¥10,000）
\item 段ボール箱と吸音材（¥2,000）
\end{itemize}
高価な装置は何も必要ありません。


\section{チェックリスト}

\begin{itemize}
\item[$\square$] OpenSCADをインストール
\item[$\square$] \texttt{pi\_berry\_crystal\_cylinders.scad} を開いてSTL生成
\item[$\square$] 対照結晶（$r_1 = r_2 = 3.5$\,mm）のSTLも生成
\item[$\square$] Bambu Lab P1S で両方を印刷（インフィル100\%）
\item[$\square$] 秋月/Amazon で測定装置を発注
\item[$\square$] 防音ボックスを作成（段ボール+吸音スポンジ）
\item[$\square$] 参照測定（結晶なし）：$T_0(f)$, $\phi_0(f)$ を記録
\item[$\square$] $\pi$結晶の透過スペクトル測定：$T(f)$, $\phi(f)$ を記録
\item[$\square$] バンドギャップ周波数を同定
\item[$\square$] バンド交差でBerry位相$\pi$の位相ジャンプを探索
\item[$\square$] 対照結晶で同じ測定 → 差分確認
\item[$\square$] 結果を論文にまとめてarXivに投稿
\item[$\square$] Phase P+ / Phase 0.5 の準備開始
\end{itemize}


\section{まとめ}

\begin{tcolorbox}[colback=green!5,colframe=green!60!black]
\textbf{Phase P の意義}

\begin{enumerate}
\item 最も安い実験（¥27,000）で、最も基本的な問いに答える：
  「Berry位相$\pi$はフォノニック結晶のバルク全体に存在するか？」

\item 成功すれば、トポロジカル・フォノニクスの新発見（独立した論文価値）。

\item 反重力への6段階のロードマップの第1段階。
  各段階はGO/STOP判定を含み、無駄なコストを避ける。

\item 全工程が室温・卓上で完結。特別な装置は不要。

\item 失敗しても（位相ジャンプが見えなくても）、
  構造最適化で再挑戦可能。理論全体が棄却されるわけではない。
\end{enumerate}

\textbf{結論}：¥27,000 と1ヶ月で、人類史上初の「反重力物質」への扉を
開く可能性のある実験を開始できる。どちらの結果も物理学に貢献する。
\end{tcolorbox}

\end{document}
